% ============================================================================
%  c_tmd snippet 2/4: 夸克中夸克的一圈 TMD（E11.15）——实辐射、虚图、软因子、组合、b 空间、
%  重整化、gamma_F、Collins--Soper 核、小 b 匹配系数；两种 MSbar 约定。
%  记号：a = alpha_s C_F/(2 pi)，L_b = ln(mu^2 b_T^2/b_0^2)，l_zeta = ln(zeta/mu^2)，
%        p(x) = (1+x^2)[1/(1-x)]_+ ，P_qq = C_F[p(x) + 3/2 delta(1-x)]。
% ============================================================================
\section{一圈计算：夸克中的夸克}
\label{sec:tmd-oneloop}

\begin{example}[一圈 TMD、软因子与 Collins--Soper 核]
\label{ex:tmd-oneloop}
\textbf{设定。}靶为在壳无质量夸克 $p=(p^+,0,\boldsymbol 0_{\mathrm T})$，对自旋与色平均；树图 $\tilde f^{(0)}=\delta(1-x)$。
记 $a\equiv\as C_F/(2\pi)$，$c_S\equiv4\pi\as\mu^{2\varepsilon}C_F/(2\pi)^{3-2\varepsilon}$（采用 Collins 的记账方式：$g_0^2=4\pi\as\mu^{2\varepsilon}$，$\MSbar$ 反项为 $S_\varepsilon\times$纯极点；它与 CONTEXT 的 $\mu^2\to\mu^2\ee^{\gamma_E}/4\pi$ 等价，重整化后结果相同），$L_b\equiv\ln(\mu^2b_{\mathrm T}^2/b_0^2)$，
$l_\zeta\equiv\ln(\zeta/\mu^2)$，$p(x)\equiv(1+x^2)[\frac1{1-x}]_+=(1-x)+2x[\frac1{1-x}]_+$。
插入单胶子末态 $l=p-k$（$l^+=(1-x)p^+$，$\boldsymbol l_{\mathrm T}=-\boldsymbol k_{\mathrm T}$）得
$f^{\mathrm{real}}=\frac{1}{4(2\pi)^{D-1}l^+}\,\frac12\sum\overline{\mathcal M}\gamma^+\mathcal M\equiv\frac{c_S}{\kT^2}F$，SIDIS 的切割振幅为
\begin{equation}
  \mathcal M^\mu=-g\,t^a\Big[\frac{\sla k\gamma^\mu}{k^2+\ii0}+\frac{n_B^\mu}{n_B\cdot l+\ii0}\Big]u(p)\qquad(\text{DY：}n_B\cdot l-\ii0),
  \label{eq:tmd-real-amp}
\end{equation}
它满足 $l_\mu\mathcal M^\mu=0$，故极化求和可取 $-g_{\mu\nu}$。

\textbf{实辐射。}取 $n_B=(s\,\ee^{2y_B},1,\boldsymbol 0_{\mathrm T})$（$s=-1$ 为 Collins 的类空线，$s=+1$ 类时线作对照），
$\delta^2\equiv\ee^{2y_B}\kT^2/(2(p^+)^2)$，$D$ 维 Dirac 迹给出（夸克--夸克、干涉、Wilson--Wilson 三类切割）
\begin{equation}
  F_a=(1-\varepsilon)(1-x),\qquad F_b=\frac{2x(1-x)}{(1-x)^2+s\delta^2},\qquad F_c=-\frac{2s\,\delta^2(1-x)}{\big((1-x)^2+s\delta^2\big)^2}.
  \label{eq:tmd-real-graphs}
\end{equation}
类空线的干涉项取主值（SIDIS 与 DY 对 $f_1$ 相同）；$F_c$ 是 Wilson 线自能，对 $s=-1$ 不可积（“无自作用”中去掉），对 $s=+1$ 有限。
对任意检验函数，$\delta\to0$ 时
\begin{equation}
  F_b\to2x\Big[\frac1{1-x}\Big]_++\delta(1-x)\ln\frac1{\delta^2},\qquad F_c^{(s=+1)}\to-\delta(1-x),
  \label{eq:rapidity-limit}
\end{equation}
（$\int_0^1F_b^{(s=+1)}\dd x=-2-2\ln\delta+O(\delta)$，$\int_0^1F_c^{(s=+1)}\dd x=-1/(1+\delta^2)$）。于是对 $\kT\neq0$
\begin{equation}
  f^{\mathrm{unsub}\,(1)}(x,\boldsymbol k_{\mathrm T};y_B)=\frac{c_S}{\kT^2}\Big\{p(x)-\varepsilon(1-x)+\delta(1-x)\Big[\ln\frac{2(p^+)^2\ee^{-2y_B}}{\kT^2}-\sigma\Big]\Big\},
  \label{eq:tmd-unsub-kspace}
\end{equation}
$\sigma=1$（类时，含自能）或 $0$（类空，无自作用）；$-2y_B$ 即快度发散。

\textbf{虚图。}外线夸克自能、Wilson 线自能以及软因子的全部虚图都是无标度积分，在维数正规化中为零。
夸克--Wilson 线顶角图却\emph{不是}无标度的：由 Lorentz 不变性、量纲与对 $n$ 的齐次性，它正比于 $(\mu^2n^2/(p\cdot n)^2)^\varepsilon$；
对类时 $n$、未来指向线，用 Feynman 参数与 Schwinger 参数可得闭式
\begin{equation}
  \begin{split}
  V^{\mathrm{SIDIS}}&=-\frac{\as C_F}{4\pi}\Big(\frac{\pi\mu^2n^2}{(p\cdot n)^2}\Big)^{\!\varepsilon}
  \frac{\Gamma(1-\varepsilon)\Gamma(1+2\varepsilon)(1-\varepsilon)}{(1-2\varepsilon)\varepsilon^2}\\
  &=-\frac{\as C_F}{4\pi}S_\varepsilon\Big[\frac1{\varepsilon^2}+\frac{1-\ell}{\varepsilon}+2+\frac{5\pi^2}{12}-\ell+\frac{\ell^2}2+\order{\varepsilon}\Big],
  \end{split}
  \label{eq:vertex-nonlightlike}
\end{equation}
$\ell=\ln(\Lambda_B^2/\mu^2)$，$\Lambda_B^2=4(p\cdot n_B)^2/n_B^2=2(p^+)^2\ee^{-2y_B}$（第二个等号用本书的 $S_\varepsilon$）。
其紫外极点 $+\frac{\as C_F}{4\pi\varepsilon}$ 与 $n$ 无关，依赖 $\ell$ 的极点是红外的。
若按 \cref{def:tmd-collins} 的次序先在固定 $\varepsilon$（$\mathrm{Re}\,\varepsilon>0$）下令 $y_B\to-\infty$，则
$|(\mu^2/\Lambda_B^2)^\varepsilon|\to0$，$V\to0$，结果再解析延拓到全部 $\varepsilon$；若先展开 $\varepsilon$，则残留 $\ell^2/2$——
两个极限不可交换。同一极限下，\eqref{eq:rapidity-limit} 的修正项集中在 $\kT\gtrsim\Lambda_B$，在固定 $b_{\mathrm T}$ 处按
$1/(\Lambda_Bb_{\mathrm T})$ 的幂次消失。\textbf{因此纯维数正规化下裸 Collins TMD 只来自实辐射与软因子的实辐射。}

\begin{remark}[类时线与类空线]
\label{rem:timelike-lines}
对过去指向（DY）线，$V$ 中出现带相位 $\ee^{\ii\pi(1+\varepsilon)}$ 的 Beta 函数。显式计算给出：对类空 $n$，
$V^{\mathrm{SIDIS}}=(V^{\mathrm{DY}})^*$，故 $2\,\mathrm{Re}\,V$ 相同，与 \cref{thm:sivers-sign} 的 PT 论证一致（类空路径上的路径编序与
时间编序重合）；对类时 $n$，$\mathrm{Re}\,V^{\mathrm{DY}}-V^{\mathrm{SIDIS}}=2\pi^2\frac{\as C_F}{4\pi}+\order{\varepsilon}$：过去指向类时线的路径编序是
反时间编序，而 Feynman 规则计算的是时间编序，二者相差入射夸克与入射程函线之间的 Coulomb 相位，在有限快度下甚至会使 $f_1$
依赖于过程。这些项都 $\propto(|n^2|)^\varepsilon$，在类光极限下消失；它们给出了 Collins\cite{Collins2011}（及
\cite{AybatRogers2011}，引 \cite{CollinsMetz2004}）采用类空正规化线的一个具体一圈理由。
\end{remark}

\textbf{软因子。}两线间的实胶子交换由色单态程函流 $J^\mu=g\,t^a[n_A^\mu/(n_A\cdot l)-n_B^\mu/(n_B\cdot l)]$ 给出（$J\cdot l=0$，故下式非负）。
记 $Y=y_A-y_B$，$l=(\tfrac{l_{\mathrm T}}{\sqrt2}\ee^{y},\tfrac{l_{\mathrm T}}{\sqrt2}\ee^{-y},\boldsymbol l_{\mathrm T})$，则
\begin{equation}
  \begin{gathered}
  -\frac{J\cdot J}{g^2C_F}=\frac{2\,n_A\cdot n_B}{(n_A\cdot l)(n_B\cdot l)}-\frac{n_A^2}{(n_A\cdot l)^2}-\frac{n_B^2}{(n_B\cdot l)^2},\\
  l_{\mathrm T}^2\frac{n_A\cdot n_B}{(n_A\cdot l)(n_B\cdot l)}=
  \begin{cases}
   \coth Y\,\big[\tanh(y-y_B)-\tanh(y-y_A)\big] & (\text{类时}),\\[1ex]
   \coth Y\,\big[\coth(y-y_B)-\coth(y-y_A)\big] & (\text{类空}),
  \end{cases}
  \end{gathered}
  \label{eq:soft-rapidity-integrand}
\end{equation}
（即类时 $\frac{\cosh Y}{\cosh(y-y_A)\cosh(y-y_B)}$、类空 $\frac{-\cosh Y}{\sinh(y-y_A)\sinh(y-y_B)}$），利用 $\int_{-L}^{L}\tanh(y-c)\,\dd y\to-2c$ 与 $\mathrm{PV}\!\int_{-L}^{L}\coth(y-c)\,\dd y\to-2c$，两种情形都得 $\int\dd y=2Y\coth Y$（类空线在
$n\cdot l=0$ 处的单极点取主值，因 $(n_A\cdot l\mp\ii0)^{-1}(n_B\cdot l\pm\ii0)^{-1}$ 的实部为两个主值之积）。类时线的两个自能项各给 $-2$；
类空线的自能 $\propto1/\sinh^2$ 不可积（线长发散，属“无自作用”）。于是 $\int\dd y\,l_{\mathrm T}^2(-J\cdot J)/(g^2C_F)=4(Y\coth Y-\sigma)$。

\begin{lemma}[一圈软因子]\tagR
\label{lem:soft-bt}
在维数正规化下（虚图无标度）
\begin{equation}
  \tilde S^{(1)}_{\text{裸}}(b_{\mathrm T};y_A,y_B)=\frac{\as C_F}{\pi}S_\varepsilon\,(Y\coth Y-\sigma)\,A(\varepsilon),\qquad
  A(\varepsilon)=-\frac1\varepsilon-L_b-\varepsilon\Big(\frac{L_b^2}2+\frac{\pi^2}{12}\Big)+\order{\varepsilon^2},
  \label{eq:soft-1loop}
\end{equation}
极点减除后 $\tilde S^{(1)}=-\frac{\as C_F}{\pi}(Y\coth Y-\sigma)L_b$。$\sigma=1$ 时极点系数是两个尖点的一圈尖点反常量纲
$\frac{\as C_F}{\pi}(Y\coth Y-1)$\cite{Polyakov1980,KorchemskyRadyushkin1987}。仅由软因子即可得
$\tilde K=\frac12\partial_{y_n}\ln\frac{\tilde S(b_{\mathrm T};y_n,-\infty)}{\tilde S(b_{\mathrm T};+\infty,y_n)}=-\frac{\as C_F}{\pi}L_b$。
\end{lemma}

\textbf{组合（动量空间）。}大快度间隔下 $\frac12[\tilde S(y_A,y_n)-\tilde S(y_A,y_B)-\tilde S(y_n,y_B)]$ 贡献
$\frac{c_S}{\kT^2}\delta(1-x)[2y_B-2y_n+\sigma]$，与 \eqref{eq:tmd-unsub-kspace} 相加，$y_A$、$y_B$ 与自能常数 $\sigma$ 全部消去：
\begin{equation}
  f^{(1)}(x,\boldsymbol k_{\mathrm T};\zeta)=\frac{c_S}{\kT^2}\Big[p(x)-\varepsilon(1-x)+\delta(1-x)\ln\frac{\zeta}{\kT^2}\Big],\qquad\zeta=2(p^+)^2\ee^{-2y_n};
  \label{eq:tmd-1loop-kspace}
\end{equation}
$\varepsilon=0$、$\kT>0$ 时即 $f^{(1)}=\frac{\as C_F}{2\pi^2\kT^2}\big[p(x)+\delta(1-x)\ln\frac{\zeta}{\kT^2}\big]$。

\textbf{横向坐标空间。}由 $\int\dd^dk\,\ee^{-\ii\boldsymbol k\cdot\boldsymbol b}(k^2)^{-\alpha}=\pi^{d/2}2^{d-2\alpha}\frac{\Gamma(d/2-\alpha)}{\Gamma(\alpha)}(b^2)^{\alpha-d/2}$，
\begin{equation}
  \begin{gathered}
  A(\varepsilon)\equiv\frac{c_S\,\mathrm{FT}[1/k^2]}{aS_\varepsilon}=-\frac{\ee^{\varepsilon L_b}\ee^{-\gamma_E\varepsilon}\Gamma(1-\varepsilon)}{\varepsilon},\\
  B(\varepsilon)\equiv\frac{c_S\,\mathrm{FT}[\ln(\zeta/k^2)/k^2]}{aS_\varepsilon}=A(\varepsilon)\Big[l_\zeta+L_b-\frac1\varepsilon-\psi(1-\varepsilon)-\gamma_E\Big],
  \end{gathered}
  \label{eq:tmd-AB}
\end{equation}
裸 Collins TMD 为
\begin{equation}
  \tilde f^{(1)}_{\text{裸}}=aS_\varepsilon\Big\{-\frac{p(x)}{\varepsilon}-L_b\,p(x)+(1-x)
  +\delta(1-x)\Big[\frac1{\varepsilon^2}-\frac{l_\zeta}{\varepsilon}-\frac{L_b^2}2-L_bl_\zeta-\kappa\frac{\pi^2}{12}\Big]\Big\}+\order{\varepsilon},
  \label{eq:tmd-1loop-bare}
\end{equation}
$\kappa=1$（本书的 $S_\varepsilon=(4\pi)^\varepsilon\ee^{-\varepsilon\gamma_E}$），$\kappa=0$（Collins 的 $S_\varepsilon=(4\pi)^\varepsilon/\Gamma(1-\varepsilon)$）。

\textbf{重整化与反常量纲。}在壳夸克靶 TMD 的红外极点必须与纯维数正规化下 $\MSbar$ 共线 PDF 的极点
$-\frac a\varepsilon[p(x)+\frac32\delta(1-x)]$ 相同（$x$ 依赖部分 $-p/\varepsilon$ 自动吻合）；其余为紫外极点：
\begin{equation}
  Z_FZ_2=1+aS_\varepsilon\Big[-\frac1{\varepsilon^2}+\frac{l_\zeta-\tfrac32}{\varepsilon}\Big],
  \label{eq:ZF-1loop}
\end{equation}
\begin{equation}
  \begin{gathered}
  \gamma_F\Big(\mu;\frac\zeta{\mu^2}\Big)=\frac{\dd\ln(Z_FZ_2)}{\dd\ln\mu}=\frac{\as C_F}{\pi}\Big(\frac32-\ln\frac\zeta{\mu^2}\Big),\\
  \tilde K(b_{\mathrm T};\mu)=\frac{\partial\ln\tilde f}{\partial\ln\sqrt\zeta}=-\frac{\as C_F}{\pi}\ln\frac{\mu^2b_{\mathrm T}^2}{b_0^2},\qquad
  \gamma_K=-\frac{\dd\tilde K}{\dd\ln\mu}=\frac{2\as C_F}{\pi}.
  \end{gathered}
  \label{eq:gammaF-K-1loop}
\end{equation}
它们满足 $\partial\gamma_F/\partial\ln\sqrt\zeta=-\gamma_K$，$\gamma_K=2\Gamma_{\mathrm{cusp}}$；$\tilde K$ 与 \cref{lem:soft-bt} 仅由软因子得到的结果相同；
与 \cite{AybatRogers2011} 式 (B1)--(B3) 一致。

\textbf{小 $b$ 匹配系数。}以重整化部分子 PDF 除重整化部分子 TMD：
\begin{align}
  \tilde C_{q/q}&=\delta(1-x)+\frac{\as C_F}{2\pi}\Big\{-L_b\,p(x)+(1-x)+\delta(1-x)\Big[-\frac{L_b^2}2-L_b\ln\frac\zeta{\mu^2}-\kappa\frac{\pi^2}{12}\Big]\Big\},
  \label{eq:Ctilde-qq}\\
  \tilde C_{q/g}&=\frac{\as T_F}{2\pi}\Big\{-L_b\,[x^2+(1-x)^2]+2x(1-x)\Big\},
  \label{eq:Ctilde-qg}
\end{align}
（$q/g$ 道：$g\to q\bar q$ 实图，胶子极化按 $D-2$ 个物理态平均，无 Wilson 线与软因子贡献。）在 $\mu=\mu_b$、$\zeta=\mu_b^2$：
$\tilde C_{q/q}=\delta(1-x)+\frac{\as C_F}{2\pi}[(1-x)-\kappa\frac{\pi^2}{12}\delta(1-x)]$，$\tilde C_{q/g}=\frac{\as T_F}{2\pi}2x(1-x)$；
$\kappa=0$ 即 \cite{AybatRogers2011} 式 (A10)、(A11)，$\kappa=1$ 即 EIS\cite{EchevarriaIdilbiScimemi2012} 第 4 节的匹配系数（$\zeta=Q^2$）。

\textbf{$\delta$ 正规子交叉检验。}未减除部分 $F_b=\frac{2x(1-x)}{(1-x)^2+r^2}$（$r=\delta^+/p^+$），软因子的快度积分 $\to\ln\frac{l_{\mathrm T}^2}{2\delta^+\delta^-}$，
\begin{equation}
  \ln\frac{(p^+)^2}{(\delta^+)^2}-\ln\frac{\kT^2}{2\delta^+\delta^-}=\ln\frac{\zeta}{\kT^2},\qquad\zeta=2(p^+)^2\frac{\delta^-}{\delta^+},
  \label{eq:delta-reg-check}
\end{equation}
与 \cite{CollinsRogers2013} 的对应完全一致（软因子虚图在此正规子下不是无标度的，但由 $\tilde S(0)=1$ 知其贡献只改变 $O(b^2\delta^+\delta^-)$）。
\end{example}

\begin{table}[htbp]
\centering
\caption{TMD 一圈有限项对 $\MSbar$ 约定的依赖（$a=\as C_F/2\pi$，$L=\ln(Q^2/\mu^2)$）。截面中的乘积 $\mathcal H\tilde C\tilde C$ 与约定无关。}
\label{tab:tmd-scheme}
\begin{tabular}{lcc}
\toprule
 & 本书 $S_\varepsilon=(4\pi)^\varepsilon\ee^{-\varepsilon\gamma_E}$ & Collins $S_\varepsilon=(4\pi)^\varepsilon/\Gamma(1-\varepsilon)$\\
\midrule
$\tilde C_{q/q}(x;\mu_b,\mu_b^2)$ & $\delta(1-x)+a[(1-x)-\frac{\pi^2}{12}\delta(1-x)]$ & $\delta(1-x)+a(1-x)$\\
$\mathcal H_{\mathrm{DY}}$ & $1+a[-L^2+3L-8+\frac{7\pi^2}6]$ & $1+a[-L^2+3L-8+\pi^2]$\\
$\mathcal H_{\mathrm{SIDIS}}$ & $1+a[-L^2+3L-8+\frac{\pi^2}6]$ & $1+a[-L^2+3L-8]$\\
$\gamma_F,\ \gamma_K,\ \tilde K$ & \multicolumn{2}{c}{一圈相同，见 \eqref{eq:gammaF-K-1loop}}\\
\bottomrule
\end{tabular}
\end{table}
